\documentclass[11pt,a4paper]{article}
\usepackage[utf8]{inputenc}
\usepackage[T1]{fontenc}
\usepackage[french]{babel}
\usepackage{lmodern}
\usepackage{amsmath,amssymb}
\usepackage[margin=2.5cm]{geometry}
\usepackage{microtype}
\usepackage{hyperref}
\usepackage{listings}
\lstset{basicstyle=\ttfamily\small, frame=single, breaklines=true,
  columns=fullflexible, keepspaces=true}
\title{Projet MLSF : résolution du jeu Meowdoku\\
  \large Question 1 : modélisation, conversion et génération de grilles}
\author{}
\date{}

\begin{document}
\maketitle

\section{Objectif et choix de conception}

L'objectif de cette question est de traduire les règles du Meowdoku en une
formule propositionnelle en forme normale conjonctive (CNF), puis de produire
un fichier DIMACS utilisable par un solveur SAT. Il faut également construire
des grilles dont l'existence d'une solution est garantie.

Nous distinguons deux programmes. Le fichier \texttt{grid\_convertion.py}
lit une grille et produit son encodage DIMACS. Cette conversion accepte le
cas général d'une grille de $N$ lignes, $M$ colonnes et $K$ régions.
Le fichier \texttt{grid\_generation.py} génère des grilles. Pour cette seconde
tâche, nous retenons la simplification autorisée par l'énoncé :
$N=M=K=n$. Le programme fonctionne avec Python 3.12 et utilise seulement
la bibliothèque standard.

\section{Modélisation propositionnelle}

\subsection{Variables et signification}

Pour chaque case $(i,j)$, avec $1\leq i\leq N$ et $1\leq j\leq M$, nous
introduisons une variable booléenne $x_{i,j}$. Elle vaut vrai si la case
contient un chat et faux sinon. Le nombre de variables est donc $NM$.
Les régions $R_1,\ldots,R_K$ constituent une partition des cases : chaque
case appartient à une seule région et chaque région est non vide.

La formule finale est la conjonction de toutes les contraintes présentées
ci-dessous. Chaque contrainte est déjà écrite sous forme de clause ou de
conjonction de clauses ; aucune transformation supplémentaire en CNF n'est
nécessaire.

\subsection{Au plus un chat par ligne}

Pour une ligne donnée, nous interdisons que deux cases distinctes contiennent
simultanément un chat. Pour chaque $i$ et chaque paire $j<\ell$, nous ajoutons :
\[
  \neg x_{i,j}\lor\neg x_{i,\ell}.
\]
Cette clause est fausse seulement lorsque les deux variables valent vrai.
Elle autorise donc zéro ou un chat parmi les deux cases. En considérant
toutes les paires de la ligne, elle garantit au plus un chat dans cette ligne.

Par exemple, pour trois cases de variables $a,b,c$, les clauses sont
$(\neg a\lor\neg b)$, $(\neg a\lor\neg c)$ et
$(\neg b\lor\neg c)$. Aucune clause positive n'est ajoutée pour imposer
un chat par ligne dans le cas général : l'énoncé demande seulement
\og au plus un \fg{}.

\subsection{Au plus un chat par colonne}

Le même principe est appliqué aux colonnes. Pour chaque colonne $j$ et
chaque paire de lignes $i<k$, nous ajoutons :
\[
  \neg x_{i,j}\lor\neg x_{k,j}.
\]
Les contraintes de lignes et de colonnes sont indépendantes des noms
et des formes des régions.

\subsection{Exactement un chat par région}

Pour chaque région $R_r$, nous combinons deux exigences. La première impose
\emph{au moins un chat} :
\[
  \bigvee_{(i,j)\in R_r}x_{i,j}.
\]
La seconde impose \emph{au plus un chat}. Pour chaque paire de cases
distinctes $u,v$ de cette région, nous ajoutons :
\[
  \neg x_u\lor\neg x_v.
\]
Ainsi, une région contenant les variables $a,b,c$ fournit la clause
$(a\lor b\lor c)$ et les trois clauses d'interdiction par paires.
La clause positive exclut une région vide de chats ; les clauses négatives
excluent une région avec plusieurs chats. Ensemble, elles imposent
exactement un chat. Pour une région réduite à une seule case, la clause
positive est unitaire et impose directement un chat sur cette case.

\subsection{Absence de chats adjacents}

Deux cases distinctes $(i,j)$ et $(k,\ell)$ sont voisines au sens du
8-voisinage lorsque
\[
  \max\bigl(|i-k|,|j-\ell|\bigr)=1.
\]
Elles ne doivent pas contenir deux chats simultanément. Les voisinages
horizontaux sont déjà couverts par les contraintes de lignes et les
voisinages verticaux par les contraintes de colonnes. Il suffit donc
d'ajouter les deux types de voisinages diagonaux :
\[
  \neg x_{i,j}\lor\neg x_{i+1,j+1},\qquad
  \neg x_{i,j+1}\lor\neg x_{i+1,j},
\]
pour $1\leq i<N$ et $1\leq j<M$. Considérer les diagonales vers la ligne
suivante suffit : chaque paire diagonale est alors traitée une seule fois.

\subsection{Correction de l'encodage}

Si une valuation satisfait la formule, nous plaçons un chat sur chaque
case dont la variable vaut vrai. Les clauses par paires garantissent
au plus un chat par ligne et par colonne ainsi que l'absence de voisinage
interdit. Les clauses des régions garantissent exactement un chat dans
chacune d'elles. Le placement respecte donc toutes les règles du jeu.

Réciproquement, un placement valide définit une valuation en mettant à
vrai les variables des cases occupées. Chaque clause positive de région
est satisfaite, puisque la région possède un chat. Chaque clause négative
est satisfaite, puisque les deux cases qu'elle concerne ne sont jamais
occupées simultanément. La formule est donc satisfiable si et seulement
si la grille possède une solution.

Dans le cas $N=M=K=n$, les contraintes de régions imposent exactement
$n$ chats. Comme il y a $n$ lignes et au plus un chat par ligne, chaque ligne
contient nécessairement un chat. Le même raisonnement s'applique aux
colonnes. Il n'est donc pas nécessaire d'ajouter des clauses positives
pour chaque ligne ou colonne dans ce cas particulier.

\section{Lecture de la grille et production du fichier DIMACS}

\subsection{Format et vérifications}

La fonction \texttt{read\_meowdoku} lit le format prescrit par le sujet.
Elle ignore les lignes vides et les commentaires commençant par
\texttt{c}. La première ligne utile indique les dimensions avec
\texttt{GRID N M}, la deuxième déclare les régions avec \texttt{ZONES},
puis chaque ligne décrit les régions auxquelles appartiennent les cases.
Par exemple :
\begin{lstlisting}
c Exemple de format
GRID 2 3
ZONES A B
A A B
A B B
\end{lstlisting}
Cet exemple illustre le format d'entrée, sans garantir que cette grille
particulière possède une solution.

Le lecteur contrôle que les dimensions sont des entiers strictement
positifs, que les noms de régions sont distincts, et que la matrice possède
le nombre annoncé de lignes et de colonnes. Il vérifie également que toutes
les régions déclarées apparaissent et qu'aucune région non déclarée n'est
utilisée. Une entrée mal formée provoque une exception \texttt{ValueError}.
Une grille bien formée peut néanmoins être insatisfiable : le lecteur
ne cherche pas de solution.

Le résultat de la lecture est une liste de listes de noms de régions.
Cette représentation permet d'accéder directement à la région d'une case
et de regrouper les cases par région dans un dictionnaire.

\subsection{Numérotation des variables}

DIMACS identifie les variables par des entiers strictement positifs.
Nous utilisons une numérotation par lignes :
\[
  \operatorname{id}(i,j)=(i-1)M+j.
\]
Pour une grille de deux lignes et trois colonnes, cela donne :
\[
\begin{matrix}
1&2&3\\
4&5&6
\end{matrix}
\]
Les indices Python commencent à zéro ; le programme utilise donc
l'expression équivalente $iM+j+1$ avec ses indices internes.

\subsection{Regroupement des paires interdites}

Une paire de numéros $(a,b)$ représente deux cases qui ne peuvent pas
contenir simultanément un chat. Elle devient la clause
$\neg x_a\lor\neg x_b$. Toutes les paires sont conservées dans un ensemble
Python, ce qui élimine les doublons. En effet, deux cases peuvent appartenir
à la même ligne et à la même région : les deux règles donnent alors
exactement la même clause.

Les numéros de chaque paire sont conservés dans l'ordre croissant.
Les paires sont triées avant l'écriture pour produire un ordre stable
des clauses. Les clauses positives de régions sont ajoutées séparément.

\subsection{Écriture DIMACS}

La fonction \texttt{meowdoku\_to\_dimacs} écrit un en-tête
\texttt{p cnf V C}, où $V=NM$ et $C$ est le nombre réel de clauses après
suppression des doublons. Chaque clause est écrite sur une ligne terminée
par \texttt{0}. Un entier positif représente une variable ; un entier
négatif représente sa négation.

Par exemple, la paire interdite $(1,2)$ donne :
\begin{lstlisting}
-1 -2 0
\end{lstlisting}
Une région contenant les cases $1$, $2$ et $4$ fournit notamment la clause
positive :
\begin{lstlisting}
1 2 4 0
\end{lstlisting}
L'exécution de la conversion peut se faire ainsi :
\begin{lstlisting}
uv run python grid_convertion.py grille.txt grille.cnf
\end{lstlisting}

\section{Génération d'une grille possédant une solution}

\subsection{Construire la solution avant les régions}

Une possibilité serait de produire des régions au hasard, puis d'appeler
un solveur et de recommencer lorsque la grille est insatisfiable.
Nous choisissons de construire d'abord un placement valide de chats,
puis de répartir les régions autour de ce placement. Cela garantit
l'existence d'une solution sans faire intervenir le solveur SAT pendant
la génération.

La fonction \texttt{generate\_meowdoku} prend une taille $n$, un chemin
de sortie facultatif et une graine aléatoire facultative. Elle retourne
une matrice de noms de régions. Le placement témoin est utilisé en interne
mais n'est pas retourné et n'est pas écrit dans le fichier de la grille.

\subsection{Placement des chats par retour arrière}

Nous plaçons un chat sur chaque ligne, en choisissant une colonne encore
libre. Une colonne utilisée est retirée des choix possibles. Cette
construction impose automatiquement un seul chat par ligne et par colonne.
Les colonnes choisies forment donc une permutation.

Pour une nouvelle ligne, nous interdisons les colonnes dont la distance
à la colonne du chat de la ligne précédente est égale à un. La colonne
identique est déjà exclue puisqu'elle est occupée. Il reste ainsi la condition
\[
  |c_{i+1}-c_i|>1,
\]
où $c_i$ désigne la colonne du chat de la ligne $i$. Vérifier seulement
la ligne précédente suffit, car les chats des lignes plus anciennes sont
éloignés d'au moins deux lignes.

Les colonnes libres sont essayées dans un ordre aléatoire. Si un choix ne
permet pas de terminer le placement, il est annulé et un autre est essayé.
Ce retour arrière évite d'abandonner toute la construction dès qu'une
décision locale conduit à une impasse.

Pour $n=4$, la permutation de colonnes $(2,4,1,3)$, avec des coordonnées
commençant à un, fournit par exemple :
\[
\begin{matrix}
\cdot&\ast&\cdot&\cdot\\
\cdot&\cdot&\cdot&\ast\\
\ast&\cdot&\cdot&\cdot\\
\cdot&\cdot&\ast&\cdot
\end{matrix}
\]
Les différences entre colonnes consécutives sont $2$, $3$ et $2$ :
aucun chat n'est adjacent à un autre.

La taille $1$ est acceptée. Les tailles $2$ et $3$ sont refusées :
pour $2$, les deux colonnes diffèrent nécessairement de un ; pour $3$,
la colonne centrale doit voisiner dans la permutation avec une colonne
dont elle diffère de un. Pour tout $n\geq4$, une solution existe :
énumérer les colonnes paires en ordre croissant, puis les colonnes impaires
en ordre croissant fournit une permutation valide. Le programme utilise
néanmoins le retour arrière aléatoire afin de varier les placements.

\subsection{Initialisation des régions et des candidats}

Chaque case contenant un chat reçoit un nom de région distinct.
Toutes les autres cases sont initialement sans région. Les noms sont
les lettres de A à Z pour $n\leq26$, puis R1, R2, etc. au-delà.

Nous enregistrons ensuite les voisins libres de chaque chat comme
\emph{candidats à la coloration}. À ce stade, ces voisins ne sont pas
encore coloriés. Une candidature contient une case et la région qui
pourrait l'occuper. Seuls les voisins horizontaux et verticaux sont
considérés pour assurer des régions connexes par les côtés.

\subsection{Croissance progressive des régions}

À chaque étape, une candidature est choisie au hasard. Si sa case est
encore vide, elle reçoit la région proposée. Nous ajoutons alors ses
propres voisins libres aux candidats. La croissance se propage donc
au-delà des voisins immédiats des chats : chaque nouvelle case coloriée
permet d'atteindre de nouvelles cases.

Une même case peut être candidate pour plusieurs régions. La première
candidature effectivement traitée fixe sa région. Les autres candidatures
pour cette case sont ensuite ignorées, car elle n'est plus vide.
Aucune case déjà attribuée n'est recoloriée.

Le tirage porte sur les candidatures, et non directement sur les cases
vides : une case proposée plusieurs fois apparaît plusieurs fois dans
la liste. Cette méthode introduit de la variété, mais ne constitue pas
un tirage uniforme parmi toutes les partitions possibles de la grille.

\end{document}
